% =====================================================================
%  2 Risikomanagement  (Aufgabenstellung 1.1: „Identifikation potenzieller
%  Projektrisiken und Gegenmaßnahmen inkl. Bewertung der Risiken nach
%  Eintrittswahrscheinlichkeit und Schwere der potenziellen Auswirkungen“)
% =====================================================================
\section{Risikomanagement}
\TODO{Kurze Hinführung}

\subsection{Methodik der Risikobewertung}
\TODO{Skalen erläutern (z. B. Eintrittswahrscheinlichkeit W und Schwere S jeweils 1--5; Risikowert R = W $\times$ S); Schwellenwerte für Handlungsbedarf; Quelle}

\subsection{Risikoliste und Gegenmaßnahmen}
\TODO{Risiken ergänzen (technisch, organisatorisch, zeitlich, fachlich); Beispielzeile ersetzen}

\begin{longtable}{@{}L{1.0cm}L{4.4cm}>{\centering\arraybackslash}p{0.6cm}>{\centering\arraybackslash}p{0.6cm}>{\centering\arraybackslash}p{0.7cm}L{6.4cm}@{}}
  \caption{Risikoliste mit Bewertung und Gegenmaßnahmen}\label{tab:risiken}\\
  \toprule
  \textbf{ID} & \textbf{Risiko} & \textbf{W} & \textbf{S} & \textbf{R} & \textbf{Gegenmaßnahme} \\
  \midrule
  \endfirsthead
  \toprule
  \textbf{ID} & \textbf{Risiko} & \textbf{W} & \textbf{S} & \textbf{R} & \textbf{Gegenmaßnahme} \\
  \midrule
  \endhead
  \bottomrule
  \endfoot
  R01 & \TODO{Beispiel: Zeitplan wird durch berufliche Belastung überschritten} & 3 & 4 & 12 & \TODO{Puffer je Meilenstein, Priorisierung nach MoSCoW} \\
  R02 & \TODO{} & & & & \TODO{} \\
\end{longtable}
\quelle{Eigene Darstellung.}

\subsection{Risikomatrix}
\TODO{Risiken in einer Matrix (W über S) einordnen; Grafik als docs/diagramme/export/risikomatrix.pdf}
\begin{figure}[H]
  \caption{Risikomatrix}\label{abb:risikomatrix}
  \centering
  \grafik[0.7\linewidth]{risikomatrix}
  \quelle{Eigene Darstellung.}
\end{figure}
